\section{Design pattern}

\subsection{Dependency injection}

\subsubsection{Descrizione}
Il \textit{pattern \textbf{Dependency Injection}} è una tecnica di progettazione \textit{software} che consente di gestire le dipendenze tra i componenti in modo più flessibile 
e modulare. Invece di creare direttamente le dipendenze all'interno dei componenti, queste vengono fornite dall'esterno (tramite "injection"). \\ \\
Sono principalmente 2 i metodi per implementare la \textit{Dependency Injection}:
\begin{itemize}
    \item \textit{\textbf{Constructor Injection}}: le dipendenze vengono fornite attraverso il costruttore del componente;
    \item \textit{\textbf{Setter Injection}}: le dipendenze vengono fornite attraverso metodi \textit{setter}.
\end{itemize}

\subsubsection{Ragione dell'utilizzo}
L'utilizzo del \textit{pattern} \textit{Dependency Injection} consente di migliorare la testabilità, la manutenibilità e la flessibilità del codice.
In particolare, permette di:
\begin{itemize}
    \item ridurre l'accoppiamento tra i componenti;
    \item centralizzare la gestione delle dipendenze, facilitando la configurazione e la manutenzione del Sistema;
    \item facilitare la sostituzione delle dipendenze con implementazioni alternative (utile durante la fase di test).
\end{itemize}

\subsection{Adapter}

\subsubsection{Descrizione}
L'\textit{\textbf{Adapter}} è un \textit{pattern} di progettazione che consente di adattare l'interfaccia di una classe esistente a quella richiesta da un \textit{client},
senza modificare il codice della classe esistente. In questo modo, è possibile integrare componenti con interfacce incompatibili senza dover riscrivere 
il codice esistente. Esistono 2 tipi di \textit{adapter}:
\begin{itemize}
    \item \textit{class adapter}: utilizza l'ereditarietà per adattare l'interfaccia di una classe esistente a quella richiesta da un \textit{client};
    \item \textit{object adapter}: utilizza la composizione per adattare l'interfaccia di una classe esistente a quella richiesta da un \textit{client}.
\end{itemize}

\subsubsection{Ragione dell'utilizzo}
L'utilizzo del \textit{pattern \textbf{Adapter}} consente di integrare componenti con interfacce incompatibili senza dover modificare il codice esistente, 
migliorando la flessibilità e la manutenibilità del Sistema. In particolare è stato utilizzato il \textit{class adapter} che ereditano le porte primarie e secondarie, 
implementando così le funzionalità richieste dalle interfacce.

\subsection{Strategy}

\subsubsection{Descrizione}
Il \textit{pattern \textbf{Strategy}} è un \textit{pattern} comportamentale che consente di definire una famiglia di algoritmi, 
incapsularli in classi separate e renderli intercambiabili. Il \textit{pattern} permette di variare l'algoritmo utilizzato 
indipendentemente dai \textit{client} che ne fanno uso, delegando la scelta della strategia concreta al contesto che la invoca. 
I componenti principali del \textit{pattern} sono:
\begin{itemize}
    \item \textit{\textbf{Strategy}}: interfaccia comune a tutte le strategie concrete, che definisce il contratto che ciascuna 
    di esse deve rispettare;
    \item \textit{\textbf{ConcreteStrategy}}: implementazione concreta dell'interfaccia \textit{Strategy}, che realizza un 
    algoritmo specifico;
    \item \textit{\textbf{Context}}: classe che mantiene un riferimento a una \textit{Strategy} e la invoca per eseguire 
    l'operazione richiesta, senza conoscere i dettagli dell'implementazione.
\end{itemize}

\subsubsection{Ragione dell'utilizzo}
Il \textit{pattern \textbf{Strategy}} è stato adottato nella sezione \textit{analytics} del Sistema per gestire le diverse 
tipologie di analisi sui dati degli impianti. Ciascuna metrica — quali il consumo energetico, la rilevazione delle presenze, 
le anomalie di impianto e la frequenza degli allarmi — è implementata come una strategia concreta che implementa 
l'interfaccia \texttt{AnalyticsStrategy}. Il servizio \texttt{AnalyticsService} funge da contesto: itera su tutte le strategie 
registrate, le esegue per l'impianto richiesto e aggrega i risultati in una lista di \textit{plot}. Questo approccio consente di:
\begin{itemize}
    \item aggiungere nuove tipologie di analytics senza modificare il codice esistente, nel rispetto del principio 
    \textit{Open/Closed};
    \item isolare ciascun algoritmo di analisi in una classe dedicata, migliorando la leggibilità e la manutenibilità del codice;
    \item testare ogni strategia in modo indipendente, semplificando la fase di verifica e validazione.
\end{itemize}

\subsection{Repository}

\subsubsection{Descrizione}
Il \textit{pattern \textbf{Repository}} è un \textit{pattern} architetturale che introduce un livello di astrazione 
tra la logica applicativa e il meccanismo di accesso ai dati. Anziché interrogare direttamente 
il \textit{database} all'interno dei servizi, le operazioni di persistenza vengono delegate a componenti 
dedicati — i \textit{repository} — che espongono un'interfaccia uniforme indipendente dai 
dettagli implementativi dello strato di persistenza. I componenti principali del \textit{pattern} sono:
\begin{itemize}
    \item \textit{\textbf{Repository Interface}}: interfaccia che definisce il contratto per le operazioni 
    di accesso ai dati, disaccoppiando la logica applicativa dall'implementazione concreta;
    \item \textit{\textbf{Repository Implementation}}: classe concreta che implementa l'interfaccia, 
    gestendo direttamente la comunicazione con il database.
\end{itemize}

\subsubsection{Ragione dell'utilizzo}
Il \textit{pattern \textbf{Repository}} è stato adottato per gestire tutte le operazioni di accesso al 
\textit{database}. Ogni \textit{repository} è composto da un'interfaccia — che funge da porta di uscita 
nell'architettura esagonale — e da un'implementazione concreta che esegue le \textit{query} SQL 
tramite un \textit{pool} di connessioni. Questo approccio consente di:
\begin{itemize}
    \item isolare completamente la logica di \textit{business} dai dettagli di accesso al \textit{database}, 
    nel rispetto del principio di separazione delle responsabilità;
    \item rendere sostituibile lo strato di persistenza senza impatto sui livelli superiori 
    dell'architettura;
    \item semplificare il testing della logica applicativa, potendo sostituire l'implementazione 
    concreta con un \textit{mock} durante la fase di verifica.
\end{itemize}

\subsection{Observer}

\subsubsection{Descrizione}
Il \textit{pattern \textbf{Observer}} è un \textit{pattern} comportamentale che definisce una relazione 
uno-a-molti tra oggetti: quando un oggetto — detto \textit{subject} o \textit{observable} — cambia 
stato, tutti gli oggetti dipendenti da esso — detti \textit{observer} o \textit{subscriber} — vengono 
notificati e aggiornati automaticamente. I componenti principali del \textit{pattern} sono:
\begin{itemize}
    \item \textit{\textbf{Observable}}: oggetto che mantiene lo stato e notifica i propri \textit{subscriber} 
    al verificarsi di un evento o di un cambiamento;
    \item \textit{\textbf{Observer/Subscriber}}: oggetto che si iscrive all'\textit{observable} per ricevere 
    le notifiche e reagire di conseguenza;
    \item \textit{\textbf{Subscription}}: rappresenta il legame tra \textit{observable} e \textit{subscriber}.
\end{itemize}

\subsubsection{Ragione dell'utilizzo}
Il \textit{pattern \textbf{Observer}} è stato adottato nel \textit{frontend Angular} tramite la libreria 
\textit{RxJS}, che ne fornisce un'implementazione matura e integrata nell'ecosistema del 
\textit{framework}. In particolare, il \textit{pattern} è stato impiegato per gestire la ricezione 
delle notifiche \textit{push} provenienti dal \textit{backend} — quali allarmi e aggiornamenti 
di stato degli impianti — e per aggiornare la UI in modo reattivo al loro arrivo, senza 
necessità di \textit{polling}. Questo approccio consente di reagire in tempo reale agli eventi di impianto, garantendo all'operatore sanitario 
una visione aggiornata della situazione senza intervento manuale.